\section*{अध्याय \ref{ch:g12:polymers} --- बहुलक}

\begin{solution}{exo:g12:polymers:1}
प्रोपीन, \ce{CH2=CH-CH3}।
\end{solution}

\begin{solution}{exo:g12:polymers:2}
\ce{-CF2-CF2-}; एक \omterm{def:g12:polymers:addition-polymer}{योगज बहुलक}: \omterm{def:g12:polymers:polymer}{एकलक} में \ce{C=C} आबंध है और
कोई अन्य उत्पाद नहीं बनता।
\end{solution}

\begin{solution}{exo:g12:polymers:3}
योगज: पॉलीएथीन, PVC, पॉलीस्टाइरीन। संघनन: PET (एक पॉलिएस्टर),
नायलॉन (एक पॉलिऐमाइड)।
\end{solution}

\begin{solution}{exo:g12:polymers:4}
सेलुलोज़, स्टार्च, लेटेक्स से बना रबड़, प्रोटीन। नायलॉन और PVC
संश्लिष्ट हैं।
\end{solution}

\begin{solution}{exo:g12:polymers:5}
हर कड़ी हर \omterm{def:g12:polymers:polymer}{एकलक} के एक समूह का उपयोग करती है; दो समूहों के साथ \omterm{def:g12:polymers:polymer}{एकलक}
दोनों ओर आबंध बना सकता है और श्रृंखला बढ़ती रहती है। एक
समूह वाला \omterm{def:g12:polymers:polymer}{एकलक} श्रृंखला को रोक देता।
\end{solution}

\begin{solution}{exo:g12:polymers:6}
$M(\ce{C2H3Cl}) = 24.0 + 3.0 + 35.5 = \qty{62.5}{g/mol}$;
$n = 125000/62.5 = 2000$।
\end{solution}

\begin{solution}{exo:g12:polymers:7}
$M(\ce{C8H8}) = \qty{104.0}{g/mol}$; $M = 2500 \times 104.0 =
\qty{2.60e5}{g/mol}$।
\end{solution}

\begin{solution}{exo:g12:polymers:8}
एक \omterm{def:g7:atoms-and-molecules:molecule}{अणु} का अम्ल समूह अगले \omterm{def:g7:atoms-and-molecules:molecule}{अणु} के \omterm{def:g11:functional-groups:alcohol}{ऐल्कोहॉल} समूह के साथ \omterm{def:g11:functional-groups:carboxylic-acid}{एस्टर} बनाता है,
जिसमें अब भी एक मुक्त अम्ल समूह है: श्रृंखला बढ़ती है। \omterm{def:g12:polymers:repeat-unit}{पुनरावृत्त
इकाई} \ce{-O-CH(CH3)-CO-}; हर कड़ी पर पानी मुक्त होता है।
\end{solution}

\begin{solution}{exo:g12:polymers:9}
थैलों की फ़िल्में शाखित श्रृंखलाओं से बनी होती हैं, जो सज नहीं सकतीं: पदार्थ
अक्रिस्टलीय, नरम और पारदर्शी होता है। दूध की बोतलें लगभग
रैखिक श्रृंखलाओं से बनी होती हैं, जो क्रिस्टलीय क्षेत्रों में सजती हैं: अधिक घनी, अधिक कठोर, और
क्रिस्टलीय क्षेत्र प्रकाश को प्रकीर्णित करते हैं।
\end{solution}

\begin{solution}{exo:g12:polymers:10}
उनकी श्रृंखलाएँ अनेक सहसंयोजी तिर्यक-बंधों से जुड़ी हैं। गरम करने से
श्रृंखलाएँ फिसल नहीं सकतीं; वह आबंध तोड़ता है और पदार्थ को पिघलाने के बजाय
नष्ट कर देता है।
\end{solution}

\begin{solution}{exo:g12:polymers:11}
$M(\ce{C12H22N2O2}) = 144.0 + 22.0 + 28.0 + 32.0 = \qty{226.0}{g/mol}$;
$n = 22600/226.0 = 100$।
\end{solution}

\begin{solution}{exo:g12:polymers:12}
\[
  \ce{H2N(CH2)6NH2 + HOOC(CH2)4COOH -> H2N(CH2)6NHCO(CH2)4COOH + H2O} .
\]
उत्पाद के एक सिरे पर अब भी एक मुक्त \omterm{def:g11:functional-groups:amine}{ऐमीन} समूह और दूसरे पर एक मुक्त अम्ल
समूह है: हर एक किसी दूसरे \omterm{def:g12:polymers:polymer}{एकलक} से अभिक्रिया कर सकता है।
\end{solution}

\begin{solution}{exo:g12:polymers:13}
प्रति \omterm{def:g12:polymers:repeat-unit}{पुनरावृत्त इकाई} पानी के दो \omterm{def:g7:atoms-and-molecules:molecule}{अणु}। पुनरावृत्त इकाइयों के $n = 1000/226.0 =
\qty{4.42}{mol}$, इसलिए पानी के \qty{8.85}{mol}:
$8.85 \times 18.0 = \qty{159}{g}$।
\end{solution}

\begin{solution}{exo:g12:polymers:14}
गरम प्राकृतिक रबड़ ऐसी श्रृंखलाओं से बना है जो एक-दूसरे पर फिसलती हैं।
सल्फ़र सेतु श्रृंखलाओं को जोड़ते हैं: वे खिंचती हैं पर वापस लौट आती हैं, और अब
बहती नहीं। तिर्यक-बंधित होने के कारण टायर को पिघलाकर नया आकार नहीं दिया जा सकता।
\end{solution}

\begin{solution}{exo:g12:polymers:15}
$M(\ce{C6H10O5}) = \qty{162.0}{g/mol}$; $n = \num{1.6e6}/162.0 = 9900$
इकाइयाँ; लंबाई $9900 \times 0.5 = \qty{4900}{nm}$, लगभग
\qty{5}{\micro m}।
\end{solution}

\begin{solution}{pb:g12:polymers:1}
\textbf{1.} एथेन-1,2-डाइऑल (दो \omterm{def:g11:functional-groups:alcohol}{ऐल्कोहॉल} समूह) और
बेंज़ीन-1,4-डाइकार्बोक्सिलिक अम्ल (दो \omterm{def:g11:functional-groups:carboxylic-acid}{कार्बोक्सिलिक अम्ल} समूह)।

\textbf{2.} एक \omterm{def:g11:functional-groups:carboxylic-acid}{एस्टर} कड़ी; पानी मुक्त होता है।

\textbf{3.} एक \omterm{def:g12:polymers:addition-polymer}{संघनन बहुलक}।

\textbf{4.} $M(\ce{C2H6O2}) = \qty{62.0}{g/mol}$; $M(\ce{C8H6O4}) =
\qty{166.0}{g/mol}$।

\textbf{5.} $62.0 + 166.0 - 2 \times 18.0 = \qty{192.0}{g/mol}$; और
$10 \times 12.0 + 8 \times 1.0 + 4 \times 16.0 = 192.0$।

\textbf{6.} $n = 38400/192.0 = 200$।

\textbf{7.} प्रति \omterm{def:g12:polymers:repeat-unit}{पुनरावृत्त इकाई} दो \omterm{def:g11:functional-groups:carboxylic-acid}{एस्टर} कड़ियाँ: लगभग 400।

\textbf{8.} उसकी श्रृंखलाएँ अधिकतर रैखिक हैं, और उनके भाग व्यवस्थित
क्षेत्रों में सीधे सजते हैं।

\textbf{9.} यह \omterm{def:g8:plastics:thermoplastic}{थर्मोप्लास्टिक} है: उसकी श्रृंखलाएँ तिर्यक-बंधित नहीं हैं और
गरम होने पर एक-दूसरे पर फिसलती हैं।

\textbf{10.} $300/0.80 = \qty{375}{g}$।

\textbf{11.} $375/25 = 15$ बोतलें।

\textbf{12.} ढक्कन, लेबल, गंदगी और कतरनें छँटाई, धुलाई और कताई के दौरान
हटा दी जाती हैं या नष्ट हो जाती हैं।

\textbf{13.} अपशिष्ट रोकना (सिद्धांत 1): इस्तेमाल की हुई वस्तु \omterm{def:g5:raw-materials-and-recycling:raw-material}{कच्चा
माल} बन जाती है।

\textbf{14.} \ce{C10H8O4 + 2H2O -> C2H6O2 + C8H6O4}।

\textbf{15.} $1000/192.0 = \qty{5.21}{mol}$।

\textbf{16.} डाइऑल: $5.21 \times 62.0 = \qty{323}{g}$; अम्ल:
$5.21 \times 166.0 = \qty{865}{g}$।

\textbf{17.} पानी के $2 \times 5.21 \times 18.0 = \qty{188}{g}$;
$1000 + 188 = 323 + 865 = \qty{1188}{g}$: द्रव्यमान संरक्षित है।

\textbf{18.} \omterm{def:g12:polymers:polymer}{एकलकों} को शुद्ध करके नए जैसा अच्छा नया PET बनाया जा सकता है,
जबकि पिघलाने से हर बार श्रृंखलाएँ थोड़ी छोटी हो जाती हैं।

\textbf{19.} एक फ़्लीस जैकेट में \qty{25}{g} की 15 बोतलें लगती हैं।
\end{solution}
